\section{Discussion}\label{sec:discussion}

We set out to evaluate the sign compiler as a scoring instrument for new item text, and the evaluation returns a specific profile of what holds, what fails, and what remains untested. We take the three contributions in turn and then state the limits that apply to all of them.

\paragraph{The instrument.} Trained at full scale, the instrument encodes the direction of its declared signs: on the trained route, every sign bit moves its own surface the intended way, although only some of those surfaces agree with the GPS across countries (Section~\ref{sec:res-bank}). The magnitude of the encoding is fragile for most dimensions. Only trust keeps its agreement with the GPS when the scenario bank is replaced, and the bank-to-bank comparison locates risk taking as the dimension that the scenario text moves most (Section~\ref{sec:res-bank}). A researcher can therefore treat the instrument's trust surface as a property of the construction that survives a change of generator, and should treat the other surfaces as properties of a particular bank until a further bank says otherwise. Naming the country does not recover the trust order that the declared signs carry, whether the signs are trained into the adapters or stated in the prompt (Section~\ref{sec:res-persona}). This is a statement about the instrument and its anchor. It is silent on whether the WVS items measure the GPS constructs in people.

\paragraph{The retrospective screen.} The screen asked whether a probe that reads only the instrument picks items whose human country differences run in the assigned direction, and the registered answer is yes, weakly. The signal is a small fraction of what two human waves achieve on the same items, its bootstrap interval reaches zero, and it is concentrated in one dimension. That dimension, risk taking, is the one the second bank does not support, and its interval does not survive an exploratory adjustment for the four dimensions tested. Trust, the dimension that does survive the second bank, does not stand apart from zero in the screen. The items assigned to positive reciprocity point the wrong way. That disagreement could arise from the instrument, from the mapping between WVS items and the GPS construct, or from limited coherence between the two surveys for this dimension, and the design cannot separate these explanations. Read as cross-survey evidence, it marks positive reciprocity as the dimension on which the instrument's readings and the human survey part ways. On the items that the adapters share with two readers of the questionnaire, the registered contrasts cannot tell the adapters apart from those readers. The readers work mostly on trust and the overlap is small, so this comparison leaves open whether the adapters add anything a careful reader would miss on the dimensions the readers do not cover.

\paragraph{The diagnostic.} The factorial shows that the declared dimensions are entangled inside the instrument. Each bit moves its own surface the intended way, but the bits also move one another's surfaces, and a single common direction carries about half of the variance across surfaces. The planned separability check fails on both routes. Prompting with the same signs is more entangled still: one common direction dominates, the trust sign moves every surface, and three own effects point the wrong way. The same pattern appears on the held-out items, where the prompted probe assigned its selected items to trust alone and did not pass the null. Training on sign-labeled pairs is thus the route that keeps the direction of each declared input. It does not deliver separate dimensions. Because the two routes carry the same declared information, the contrast describes two instruments, and it makes no claim about which training recipe is best or about screening accuracy, where the difference between routes is unresolved.

\paragraph{Fidelity to labels and agreement with respondents.} The descriptive follow-up analyses (Sections~\ref{sec:res-human} to \ref{sec:res-recovery}) separate three parts of the construction that are easy to run together: the instrument, which turns declared signs into training labels; the anchor, which supplies the signs; and the model, which carries them. Fidelity to the labels is high. On training scenarios they never saw, the adapters' margins exceed the base model's on \cl{0.940}{C9_label_recovery_recovery_profile_block} of held-out pairs, and the adapters prefer their own chosen completion on \cl{0.816}{C9_label_recovery_accuracy_profile_block} of them, against \cl{0.489}{C9_label_recovery_base_accuracy_profile_block} for the base model. Agreement with respondents is a different fact. The same adapters sit farther from the survey's answer distributions than the base model, by \cl{0.045}{C9_TVD_adapter_minus_base_profile_block} \ci{0.006}{0.093}{C9_TVD_adapter_minus_base_profile_block} in total variation distance, whereas a country-named prompt, which involves no training, sits closer, by \cl{-0.105}{C9_TVD_persona_minus_base_profile_block} \ci{-0.122}{-0.086}{C9_TVD_persona_minus_base_profile_block}. The adapters' order of countries agrees with the GPS on trust and does not detectably agree with the human composite on any of seven surfaces. A model can therefore carry its declared inputs faithfully and still lie farther from respondents than a prompt that was never trained on those labels, and a construction check cannot stand in for a check against people. We read this as a point about how such an instrument should be evaluated, with the instrument, the anchor and the model each needing its own check, and not as a result about any one of them. The evidence is descriptive, it was not preregistered, it rests on one anchor and one human source, and it earns no claim. We do not test why the distributions move apart. The sign labels carry no information about magnitudes, the readout is a scoring convention, and the base model is shared, and the design cannot separate these.

\paragraph{Limits.} Five limits apply to every result above. First, the human criterion is retrospective. The WVS Wave~7 file existed before the probe, the pool and the adapters, and it is the only human source. Registration and the freeze guard protect against tuning the probe to the criterion, but they cannot make the criterion independent of the survey program from which the item pool was drawn. Pre-piloting, in the sense that screening with the instrument would change which items a researcher fields relative to a human-only rule declared in advance, is a prospective use that this paper does not test. Second, the instrument is built from signs. It carries no information about magnitudes, and it cannot represent differences among countries that share a sign profile or among people within a country. Its readings are group-level, and country-level agreement does not establish individual-level validity or measurement invariance \citep{davidov2014,bratt2025}. Third, every adapter shares one base model, one training recipe and one readout convention, and the readout's code probabilities are a scoring convention whose dispersion need not match human response distributions. Fourth, the per-dimension cells are small and the selected set churns under readout variants, so the instrument's item-level assignments are more stable than any particular list of selected items. Fifth, the noise checks are narrow: one unseeded retrain bounds retraining noise for one unit, and two near-zero refits are a screen with no claim across countries.

\paragraph{Intended use.} Within those limits, the instrument has a modest and inspectable use. A survey designer can read the dimension and direction it assigns to candidate items as one input beside expert review and human pretesting, keeping in view that no dimension passed both checks. Its value lies in what it makes explicit: the declared moments, the sign rule, and the absence of any group name at evaluation are all on the record, so a reader can see what information the instrument was given and check where its readings depart from people. It complements fielding and pretesting, and the evidence here gives no reason to use it in their place. In the terms of \citet{broska2025}, its readings are potentially informative observations that still require a human criterion.
